\documentclass{article}
\usepackage{amsmath}
\usepackage{amssymb}
\usepackage{amsthm}
\usepackage{amsfonts}
\usepackage{amscd}

\title{Research Statement}
\author{Alireza Noroozi}
% \email{anoroozi25@ku.edu.tr}
\date{September 2026}

\begin{document}

\maketitle

% Unsupervised Joint Representation Learning for Protein Sequence and Structure, with a Focus on Binder Design

% Lead supervisor: Kiarash Jamali (Principal Investigator, Generative Biology Institute, EIT) - kiarash.jamali@eit.org

% Co-supervisor: Charlotte Deane (Professor of Structural Bioinformatics, Department of Statistics, University of Oxford) - deane@stats.ox.ac.uk

Protein and binder design are central to therapeutic discovery.
A binder that recognises a chosen epitope with high affinity and selectivity can inhibit a pathway,
recruit an E3 ligase, or serve as a modular starting point for drugs and diagnostics.
Over the last few years, structure-prediction and generative models---AlphaFold-class predictors, protein language models,
and diffusion-based designers such as RFdiffusion---have been developed largely to make this problem tractable.
These methods have produced impressive successes on well-studied, globular targets. 
They remain, however, far from a general solution: several systematic failure modes
still limit their use on the proteins that matter most in disease.

The first of these failures is generalisation. Current models are trained on a relatively small
and redundant set of experimental complexes, so they behave poorly on new targets with low sequence
or structural homology, where evolutionary couplings and PDB-like templates are weak or absent.
Even when a complex can be generated, binder residue specificity is often wrong: the interface may
look geometrically plausible, yet it does not encode a chemically precise contact map,
and affinity is driven by a few non-selective interactions rather than a defined epitope.
A related and more severe problem appears on conformationally dynamic proteins. Targets such as
RBX1 undergo large motions and contain intrinsically disordered regions; they are poorly represented
by a single static backbone. Designs that assume a rigid pocket therefore bind weakly,
engage the wrong conformation, or fail altogether.

A further limitation is that many pipelines reward stickiness rather than true molecular recognition.
Models trained on sparse complex structures tend to produce hydrophobic, promiscuous interfaces that
score well in silico but exhibit cross-reactivity in cells. Distinguishing a selective binder from a
generic sticky protein remains an open problem, especially when experimental negatives and off-target complexes 
are scarce. Current generators are also restricted in chemical modality: most sample only canonical amino-acid 
sequences and standard protein backbones, even though non-canonical amino acids and peptides can expand the accessible
chemistry, improve protease resistance, and create contacts that natural residues cannot form. Without representations 
that jointly embed sequence, structure, and a broader chemical alphabet, design remains confined to a narrow and already 
well-sampled region of binder space.

Evaluation compounds these modelling gaps.
Binder design is still judged mainly with static, non-visual metrics---interface pLDDT,
predicted aligned error, Rosetta-style energy terms, or a single predicted complex---which neither 
capture induced fit nor show whether the binder remains specific across a conformational ensemble.
Such scores correlate only weakly with wet-lab success and give little insight into why a design failed.
Finally, generated candidate sets are often redundant: diffusion and inverse-folding models collapse onto
a few high-likelihood motifs, so large in silico libraries contain little true sequence or structural diversity,
and experimental capacity is spent retesting near-duplicates rather than covering distinct binding modes. 
Together, these limitations---poor behaviour on distant targets, weak residue-level specificity, difficulty
with mobile and disordered proteins, cross-reactive sticky designs, a narrow set of chemical modalities,
static evaluation, and low library diversity---define the gap that a joint sequence--structure 
representation for binder design must address.  

\end{document}